\documentclass[12pt]{article}
\textheight=24.5cm \textwidth=16.5cm
\parskip = 0.165cm
\renewcommand\baselinestretch{1.2}
\topmargin=-2cm \oddsidemargin=0cm \evensidemargin=0cm

\usepackage{amsfonts}
\usepackage{amsmath}
\usepackage{amssymb}
\usepackage{amscd}
\usepackage{latexsym}
\usepackage[round, sort]{natbib}
\usepackage[colorlinks=true,linkcolor=black,anchorcolor=black,citecolor=black,urlcolor=black]{hyperref}
\usepackage{color}
\usepackage{xcolor}
\usepackage{array}
\usepackage{colortbl}
\usepackage{algorithmicx}
\usepackage{algorithm}
\usepackage{algpseudocode}
\usepackage{enumerate}
\usepackage{float}
\usepackage{indentfirst}
\usepackage{CJKutf8}

\newtheorem{theorem}{Theorem}[section]
\newtheorem{cor}{Corollary}[section]
\newtheorem{lemma}{Lemma}[section]
\newtheorem{rem}{Remark}[section]
\newtheorem{definition}{Definition}[section]
\newtheorem{obs}{Observation}
\newtheorem{proposition}[theorem]{Proposition}
\newtheorem{exam}{Example}[section]

\title{\LARGE\bf Two-Machine Flow Shop Scheduling to Minimize Total Weighted Completion Time: Structural Properties and Dynamic Programming}
\author{Author One$^{a,}$\footnote{Corresponding author. E-mail address: author@email.com}, Author Two$^{b}$\\
{\small $^{a}$Institution One, Location, Country}\\
{\small $^{b}$Institution Two, Location, Country}\\}
\date{}

\begin{document}
\begin{sloppypar}

\maketitle

\begin{abstract}
    Minimizing the total weighted completion time on a two-machine flow shop, 
    denoted $F2 \parallel \sum w_j C_j$, is a classical NP-hard scheduling problem with applications in manufacturing, computing, and logistics. 
    Existing exact methods rely on enumerative techniques that scale poorly with instance size,
     while approximation algorithms often impose restrictive assumptions on processing times.
      Here we partition the job set into at most $K$ categories by non-increasing weight, which decomposes the schedule into $K$ sequential blocks. 
      We establish four structural properties of optimal schedules: the permutation property, the positional constraint $\pi(J_k) \le k$, 
      the block-end condition requiring the final job of each block to carry that block's weight class, and the SPT(1)-LPT(2) (Johnson) ordering rule within each block. 
      Building on these properties and the interval-based dynamic programming framework developed for $F2|r_j|C_{\max}$, 
      we design a pseudo-polynomial time dynamic programming algorithm that constructs the schedule block by block.
      We further show that geometric rounding of weights and processing times converts this algorithm into a polynomial-time approximation scheme.
\end{abstract}

\noindent \textbf{Keywords:} Flow shop scheduling; Total weighted completion time; Dynamic programming; Structural properties

\section{Introduction}
\label{sec:intro}
% 介绍两台机器流水车间问题背景、重要性
% 现有文献综述 (Literature Review)
% 本文的主要贡献与论文结构编排

\section{Problem Formulation}
\label{sec:problem}
% 确立数学模型：给出 F2 || sum W_j C_j 的问题定义，给出参数如 A_j, B_j, W_j 的定义
% 进行工件集的预处理（如按权重分至多K类等）

\section{Structural Properties}
\label{sec:properties}

We partition the job set $J$ into $K$ weight categories $J_1, J_2, \dots, J_K$ such that all jobs in $J_k$ share the same weight $W_k$, 
with $W_1 > W_2 > \dots > W_K$. Let $n_k = |J_k|$ and $\sum_{k=1}^K n_k = n$.

The schedule is divided into $K$ contiguous logical segments (blocks) $S_1, S_2, \dots, S_K$, 
processed sequentially on both machines. We denote by $\pi(j)$ the block index to which job $j$ is assigned.

\begin{lemma}
    \label{lem:permutation}
    \textbf{(Permutation property.)} For $F2 \parallel \sum w_j C_j$, 
    there exists an optimal schedule in which the job sequence is identical on both machines.
\end{lemma}

This is a classic result for two-machine flow shops with regular objective functions.

\begin{lemma}
    \label{lem:positional}
    \textbf{(Positional constraint.)} In any optimal schedule, $\pi(J_k) \le k$ for all $k = 1, \dots, K$.
     Equivalently, every job from weight class $J_k$ is assigned to one of the first $k$ blocks $S_1, S_2, \dots, S_k$.
\end{lemma}

Lemma~\ref{lem:positional} implies that once block $S_k$ has been processed, all jobs carrying weights $W_1, \dots, W_k$ are fully scheduled. 
Consequently, block $S_k$ may contain only jobs from classes $J_k, J_{k+1}, \dots, J_K$.

\begin{lemma}
    \label{lem:block_end}
    \textbf{(Block-end condition.)} In any optimal schedule, the final job of block $S_k$ belongs to weight class $J_k$, for each $k = 1, \dots, K$.
\end{lemma}

\begin{lemma}
    \label{lem:spt_lpt}
    \textbf{(Intra-block ordering.)} Within each block $S_k$, all jobs are sequenced according to the SPT(1)-LPT(2) rule (Johnson's rule): 
    for any two jobs $x, y \in S_k$, job $x$ precedes job $y$ if and only if
    \[
        \min\{a_x, b_y\} < \min\{a_y, b_x\},
    \]
    with ties broken by non-decreasing $a_j$.
\end{lemma}

Taken together, these four properties determine the structure of any optimal schedule: 
jobs are organised into $K$ blocks ordered by decreasing weight; 
within each block the sequence follows Johnson's rule; class $k$ jobs cannot appear beyond block $k$; 
and the final position of block $k$ is occupied by a job of class $k$.

\section{Dynamic Programming Algorithm}
\label{sec:dp_algorithm}

The structural properties established in Section~\ref{sec:properties} enable a dynamic programming algorithm that builds the schedule block by block. 
Our approach adapts the interval-based state representation of , originally devised for the $F2|r_j|C_{\max}$ PTAS, 
to the weighted completion time objective. In that framework, the schedule is partitioned into intervals, 
each interval is internally sequenced by Johnson's rule, and the DP state tracks per-interval completion times. 
We retain the block-structured state space but replace the makespan accumulation with an incremental weighted sum computation.

\noindent\textbf{Notation.}
\begin{CJK}{UTF8}{gbsn}
{\small\renewcommand{\arraystretch}{0.9}
\begin{tabular}{@{}c|l@{}}
    $J$        & 全体工件集合 \\
    $n$        & 工件总数 ($n = |J|$) \\
    $K$        & 不同权重的类别数 \\
    $J_k$      & 权重为 $W_k$ 的工件集合 ($k=1,\dots,K$) \\
    $W_k$      & $J_k$ 中工件的公共权重，满足 $W_1 > W_2 > \dots > W_K$ \\
    $n_k$      & $|J_k|$，第 $k$ 类工件的数量 \\
    $a_j$      & 工件 $j$ 在机器 1 上的加工时间 \\
    $b_j$      & 工件 $j$ 在机器 2 上的加工时间 \\
    $w_j$      & 工件 $j$ 的权重 ($j \in J_k \iff w_j = W_k$) \\
    $L_k$      & $J_k$ 中工件按 Johnson 规则排序后的有序列表 \\
    $S_k$      & 调度中第 $k$ 个连续块 \\
    $J_{i,k}$  & 第 $k$ 个块中的第 $i$ 个工件 ($i = 1,\dots,|S_k|$) \\
    $\mathcal{Y}^{(k)}$ & 处理完前 $k$ 个块后的非支配状态集合 ($k=0,\dots,K$) \\
    $A$                  & 前 $k$ 个块在机器 1 上累计加工时间 \\
    $B$                  & 第 $k$ 个块在机器 2 上的完工时间（即第 $k{+}1$ 个块开始时机器 2 可用时刻） \\
    $z$                  & 前 $k$ 个块的累计加权完工时间和 \\
    $A_{\text{cur}}, B_{\text{cur}}, z_{\text{cur}}$ & 块构建过程中机器 1 累计时间、机器 2 完工时间、加权成本的运行值 \\
    $Z^*$      & 最优（最小）总加权完工时间 \\
\end{tabular}}
\end{CJK}
\smallskip

\subsection{Preprocessing}

We first partition the $n$ jobs into $K$ sets $J_1, J_2, \dots, J_K$ by weight, with $W_1 > W_2 > \dots > W_K$. 
Within each set $J_k$, we sort the jobs according to Johnson's rule (SPT(1)-LPT(2)). Let the resulting ordered list be
\[
    L_k = \bigl(j_{k,1}, j_{k,2}, \dots, j_{k,n_k}\bigr), \quad k = 1,2,\dots,K,
\]
where $n_k = |J_k|$. Since block $S_k$ will contain jobs drawn from classes $k, k+1, \dots, K$ and
must obey Johnson ordering internally (Lemma~\ref{lem:spt_lpt}),
the correct within-block sequence is obtained by taking prefixes of each $L_m$ ($m \ge k$) and merging them under the same Johnson comparator.
To merge a set of jobs into Johnson order, one partitions it into those with $a_j \le b_j$ (sorted by non-decreasing $a_j$) and those with $a_j > b_j$ (sorted by non-increasing $b_j$), and concatenates the two subsequences. This procedure is used whenever the algorithm merges several $L_m$ suffixes into a single sequence.

\subsection{Dynamic Programming Algorithm}

We build the schedule block by block. After the first $k$ blocks are fixed, let
\[
    S_m = (J_{1,m}, J_{2,m}, \dots, J_{|S_m|,m}), \quad m = 1,\dots,k,
\]
where $J_{i,m}$ denotes the $i$-th job in block $S_m$. Each pre-sorted list $L_m$
is the Johnson order of class $J_m$. Since blocks consume prefixes of each $L_m$
(Lemma~\ref{lem:spt_lpt}), the jobs already assigned to $S_1 \cup \dots \cup S_k$
constitute a prefix of each $L_m$. Consequently, the set of remaining (unassigned)
jobs of class $J_m$ is simply the suffix of $L_m$ after that prefix.

A state in $\mathcal{Y}^{(k)}$ is a tuple $(J_{i,k},\, A,\, B,\, z)$, where $J_{i,k}$
collectively denotes the jobs assigned to each block $S_1,\dots,S_k$ (i.e.\
$S_m = (J_{1,m},\dots,J_{|S_m|,m})$ for $m \le k$), $A$ is the total M1 processing
time over blocks $1..k$, $B$ is the M2 completion time at the end of block $k$,
and $z$ is the accumulated weighted completion time. Since blocks consume prefixes
of each $L_m$ (Lemma~\ref{lem:spt_lpt}), the jobs assigned to $S_1 \cup \dots \cup S_k$
constitute a prefix of each $L_m$; consequently, the set of remaining (unassigned)
jobs of class $J_m$ is simply the suffix of $L_m$ after that prefix. The initial
state is $\mathcal{Y}^{(0)} = \{(\emptyset,\, 0,\, 0,\, 0)\}$ with all blocks empty.

To construct block $S_{k+1}$ from a state in $\mathcal{Y}^{(k)}$, we collect the
remaining suffixes of $L_{k+1},\dots,L_K$ and merge them into a single
Johnson-ordered sequence. Scanning this sequence from left to right, for each
job $j \in \bigcup_{m=k+1}^K J_m$ we decide whether to \textbf{include} it in
$S_{k+1}$ or \textbf{postpone} it to a later block:

\begin{itemize}
    \item If $j \in J_{k+1}$: \textbf{must include} (Lemma~\ref{lem:positional}).
    \item If $j \in J_m$ with $m > k{+}1$: either include or postpone (binary branch).
\end{itemize}

When a job is included, it becomes the next $J_{i,k+1}$ and we update:
\begin{align}
    A_{\text{cur}} &\gets A_{\text{cur}} + a_j, \notag\\
    B_{\text{cur}} &\gets \max(B_{\text{cur}}, A_{\text{cur}}) + b_j, \notag\\
    z_{\text{cur}} &\gets z_{\text{cur}} + w_j \cdot B_{\text{cur}}. \notag
\end{align}
A postponed job remains in its suffix and participates in the next block's scan.
At the end, we require that the last job of $S_{k+1}$ belongs to $J_{k+1}$
(Lemma~\ref{lem:block_end}); otherwise the branch is discarded.

A valid branch produces a successor state $(J_{i,k+1},\, A_{\text{cur}},\, B_{\text{cur}},\, z_{\text{cur}})$
where $J_{i,k+1}$ extends the previous $J_{i,k}$ with the newly formed $S_{k+1} = U$.
It is inserted into $\mathcal{Y}^{(k+1)}$ under the dominance rule:
among states with the same block partition $(S_1,\dots,S_{k+1})$, only the one
with minimal $z_{\text{cur}}$ for each $(A_{\text{cur}}, B_{\text{cur}})$ is retained.

The complete procedure is given in Algorithm~\ref{alg:dp}.

\begin{algorithm}[H]
    \renewcommand{\baselinestretch}{1.2}\selectfont
    \caption{DP\_F2\_WeightedSum}
    \label{alg:dp}
    \begin{algorithmic}[1]
        \Require Jobs partitioned into $J_1,\dots,J_K$ with $W_1>\dots>W_K$; each $L_k$ is the Johnson-ordered list of $J_k$.
        \Ensure optimal weighted sum $Z^*$

        \State $\mathcal{Y}^{(0)} \gets \{ (\emptyset,\, 0,\, 0,\, 0) \}$ \Comment{no block built yet}

        \For{$k = 0$ to $K-1$}
            \State $\mathcal{Y}^{(k+1)} \gets \emptyset$
            \For{each $(J_{i,k},\, A,\, B,\, z) \in \mathcal{Y}^{(k)}$}
                \State Merge the remaining suffixes of $L_{k+1},\dots,L_K$ in Johnson order into a sequence $\sigma$
                \State Initialize a stack with $(\ell,\, A_{\text{cur}},\, B_{\text{cur}},\, z_{\text{cur}},\, U,\, \textit{last}) \gets (1,\, A,\, B,\, z,\, \emptyset,\, \text{nil})$
                \While{stack is not empty}
                    \State Pop $(\ell,\, A_{\text{cur}},\, B_{\text{cur}},\, z_{\text{cur}},\, U,\, \textit{last})$
                    \If{$\ell > |\sigma|$}
                        \If{$U \neq \emptyset$ \textbf{and} $\textit{last} = k{+}1$}
                            \State Let $S_{k+1} \gets U$; $J_{i,k+1} \gets J_{i,k} \cup \{S_{k+1}\}$
                            \State Advance the taken-prefix pointer of each $L_m$ by $|U \cap J_m|$
                            \If{$\nexists\,(J_{i,k+1},\, A_{\text{cur}},\, B_{\text{cur}},\, \hat{z}) \in \mathcal{Y}^{(k+1)}$ with $\hat{z} \le z_{\text{cur}}$}
                                \State Insert $(J_{i,k+1},\, A_{\text{cur}},\, B_{\text{cur}},\, z_{\text{cur}})$ into $\mathcal{Y}^{(k+1)}$, pruning dominated states
                            \EndIf
                        \EndIf
                        \State \textbf{continue}
                    \EndIf
                    \State $j \gets \sigma[\ell]$, let $m$ be the class index of $j$ 
                    \If{$m = k{+}1$}
                        \State $A' \gets A_{\text{cur}} + a_j$;\; $B' \gets \max(B_{\text{cur}}, A') + b_j$;\; $z' \gets z_{\text{cur}} + w_j \cdot B'$
                        \State Push $(\ell{+}1,\, A',\, B',\, z',\, U \cup \{j\},\, k{+}1)$ onto stack
                    \Else
                        \State $A' \gets A_{\text{cur}} + a_j$;\; $B' \gets \max(B_{\text{cur}}, A') + b_j$;\; $z' \gets z_{\text{cur}} + w_j \cdot B'$
                        \State Push $(\ell{+}1,\, A',\, B',\, z',\, U \cup \{j\},\, \textit{last})$ onto stack 
                        \State Push $(\ell{+}1,\, A_{\text{cur}},\, B_{\text{cur}},\, z_{\text{cur}},\, U,\, \textit{last})$ onto stack 
                    \EndIf
                \EndWhile
            \EndFor
        \EndFor
        \State $Z^* \gets \min\{\, z \mid (J_{i,K},\, A,\, B,\, z) \in \mathcal{Y}^{(K)} \,\}$
        \State \Return $Z^*$
    \end{algorithmic}
\end{algorithm}

\section{Conclusion}
\label{sec:conclusion}
% 结论及未来展望

\bibliographystyle{apalike}
\bibliography{ref_IJPR}

\end{sloppypar}
\end{document}
